\documentclass[11pt]{article}
\usepackage[T1]{fontenc}
\usepackage[utf8]{inputenc}
\usepackage{amsmath,amssymb,amsthm,booktabs}
\usepackage[margin=1in]{geometry}
\usepackage{hyperref,xurl}
\newtheorem{theorem}{Theorem}
\newtheorem{lemma}[theorem]{Lemma}
\newtheorem{proposition}[theorem]{Proposition}
\newcommand{\HH}{\mathbb H}
\newcommand{\Q}{\mathbb Q}
\title{A triangle incidence constraint at an exceptional\linebreak compact hyperbolic Coxeter vertex}
\author{Coxeter Polytopes Frontier Project}
\date{18 October 2026}
\begin{document}
\maketitle
\begin{abstract}
At a finite vertex of type $E_7^3\sqcup E_8$ in a compact hyperbolic
Coxeter polytope in $\HH^{29}$, we designate three two-face angles in
each $E_7$ component. At least six of these nine faces are nontriangular.
The proof reduces the actual third side of a triangle to a complete
finite contact domain, checked by two exact implementations. All three
angles force the same contact vector. The common finite vertex fixes
the signs of its Lorentzian heights, and two such third sides in different
components have a product strictly between consecutive native Coxeter
bonds. No arithmeticity assumption is used. This is a local incidence
theorem, not a new dimension bound or a face-charge saving.
\end{abstract}

\section{Statement and geometric conventions}
Wall normals have square $2$ in a real Lorentz space of signature $(n,1)$.
A chamber is the intersection of their nonpositive halfspaces in the future
hyperboloid. Products of distinct outward normals are nonpositive; a finite
bond of label $m$ has product $-2\cos(\pi/m)$. Parallel and ultraparallel
walls have products $-2$ and less than $-2$, respectively.

We use the standard simplicity of compact Coxeter polytopes and the
classification of finite Coxeter diagrams. In particular, a finite vertex
of a compact $n$-polytope has exactly $n$ independent incident walls.
A whole $E_8$ has no connected finite extension, a connected finite
one-root extension of $E_7$ is $E_8$, and a finite bond of label at least
$6$ is an isolated dihedral factor. These standard inputs are separated
from the new geometric argument and finite computation below.

Number each $E_7$ by the simple edges
\[
0-1-2-3-4-5,\qquad 2-6.
\]
At a simple vertex, a pair of incident facets specifies the two-face whose
supporting walls are all the other vertex facets.

\begin{theorem}\label{thm:main}
Let $v$ be a vertex of a compact hyperbolic Coxeter polytope in
$\HH^{29}$ with wall diagram $E_7\sqcup E_7\sqcup E_7\sqcup E_8$.
In each $E_7$ designate the pairs $(2,4),(2,5),(3,5)$. Among the three
$E_7$ components, at most one supplies a triangular two-face at any of
its designated pairs. Thus at most three of the nine designated faces
are triangles, and at least six have four or more sides.
\end{theorem}

\section{The actual third side has no outside contacts}
Let $T$ be the positive span of the $29$ roots at $v$. Suppose the two-face
of a designated pair $a,b$ is triangular, with third wall root $c$.
Its other vertices have the old wall sets with $a$, respectively $b$,
replaced by $c$. These are actual positive-definite corners, not merely
subdiagrams of a chosen local witness.

\begin{lemma}\label{lem:outside}
The root $c$ is orthogonal to all old components except the marked $E_7$.
\end{lemma}
\begin{proof}
The untouched $E_8$ appears whole at both new corners, so finite maximality
forces orthogonality to it. If $c$ contacts an untouched $E_7$ component
$F$, then $F+c$ is a whole $E_8$ at each new corner. Finite maximality
therefore forces $c$ to be orthogonal to every other old root in those
corners. Their union contains every old root. Hence the Gram matrix on
all old roots and $c$ is the positive-definite direct sum
$(F+c)\sqcup(T-F)$, of rank $30$. A Lorentz space of signature $(29,1)$
has positive index $29$, a contradiction.
\end{proof}

\begin{lemma}\label{lem:labels}
Every contact label of $c$ with the marked $E_7$ belongs to
$\{2,3,4,5\}$.
\end{lemma}
\begin{proof}
Every contact occurs at one of the two finite corners, so parallel and
dotted products are excluded. A bond at least $6$ would be an isolated
$I_2(m)$ component. If its old endpoint $j$ is white, it occurs in both
corners. Its old neighbors would all have to lie in both omitted
singletons $\{a\}$ and $\{b\}$, impossible in the connected $E_7$.
If its old endpoint is a mark, that mark must be a leaf whose only
neighbor is the other, omitted mark. The designated pairs are all
nonadjacent, ruling this out.
\end{proof}

\section{Complete exact triangle domains}
Let $G$ be the $E_7$ Gram matrix and $J=G^{-1}$. Write $b$ for the seven
products of $c$ with that component. If $e$ is the future unit vector
at $v$, Lemma~\ref{lem:outside} gives
\[
c=p+h e,\qquad p\in T,\quad e^2=-1,\quad
h=-c\cdot e>0,\qquad h^2=bJb^T-2.
\]
The strict inequality holds because the third side does not pass through
$v$. Thus a nonpositive computed $h^2$ cannot supply this side.

\begin{lemma}\label{lem:affine}
The wall set of a compact acute hyperbolic chamber contains no
positive-semidefinite singular principal Coxeter diagram.
\end{lemma}
\begin{proof}
Choose a connected singular component and positive kernel coefficients
$\lambda_i$. Then $u=\sum_i\lambda_i r_i$ is null. A strict interior
future timelike point $x$ has $u\cdot x<0$, so $u$ is nonzero and future.
Its products with the component vanish, and all its other wall products
are nonpositive by acuteness. The resulting ideal point belongs to the
closed chamber, contrary to compactness. This proof does not assume that
the singular principal diagram was already known to support a face.
\end{proof}

\begin{proposition}\label{prop:finite}
For each designated pair, the unique possible triangle contact vector is
\[
c\cdot r_5=-2\cos(\pi/5),\qquad c\cdot r_j=0\quad(j\ne5).
\]
Its height satisfies $h^2=(1+3\sqrt5)/4$.
\end{proposition}
\begin{proof}[Exact finite certificate]
By Lemma~\ref{lem:labels}, enumerate all $4^7$ vectors for each pair.
Test the two rank-seven finite corners obtained by omitting each mark.
For every eligible vector compute $bJb^T-2$. Discard nonpositive values;
otherwise test the principal subsets of the eight-root diagram for the
obstruction in Lemma~\ref{lem:affine}. Exact results are:
\begin{center}
\begin{tabular}{lrrrrr}
\toprule
Pair & Vectors & Finite eligible & $h^2\le0$ & Singular PSD & Remaining\\
\midrule
$(2,4)$ & 16,384 & 12 & 3 & 8 & 1\\
$(2,5)$ & 16,384 & 7 & 3 & 3 & 1\\
$(3,5)$ & 16,384 & 7 & 3 & 3 & 1\\
\bottomrule
\end{tabular}
\end{center}
In all three rows the remaining vector is the one stated. The certificate
contains every eligible vector and each explicit singular principal
witness. The independent checker reconstructs the entire $49,152$-vector
finite filter, recomputes the heights, verifies the singular witnesses,
and exhausts all $255$ nonempty principal subsets of each survivor.

All arithmetic takes place in $\Q(\sqrt2,\sqrt5)$ because the preceding
geometric reduction forced labels at most $5$. This is an exact
computational representation, not an arithmeticity hypothesis. The
producer's finite-diagram classifier and the checker's canonical rooted
tree catalogue are separate implementations; their domain digests agree.
Exact PSD tests use different Schur elimination orders. The old local
witness catalogue also contains these three triangle recipes, but their
sufficiency is not required for the necessary classification here.
\end{proof}

\section{Coupling at the common vertex}
\begin{proof}[Proof of Theorem~\ref{thm:main}]
Suppose different $E_7$ components supply designated triangles with
third roots $c_i=p_i+h e$ and $c_j=p_j+h e$. Their projections are
orthogonal by Lemma~\ref{lem:outside}. Both heights have the same positive
sign, because the common vertex lies strictly on the interior side of
both new walls. Their products are therefore
\[
c_i\cdot c_j=-h^2,\qquad
(c_i\cdot c_j)^2=\frac{23+3\sqrt5}{8}.
\]
The roots are distinct, having nonzero contacts on different old factors.
Directed rational cosine enclosures give
\[
\bigl(2\cos(\pi/11)\bigr)^2
<\frac{37}{10}<\frac{23+3\sqrt5}{8}
<\frac{93}{25}<\bigl(2\cos(\pi/12)\bigr)^2.
\]
The middle value is less than $4$ and lies strictly between consecutive
finite-bond squares. It is neither a finite Coxeter product nor a parallel
or dotted product. Such two wall roots cannot coexist. At most one $E_7$
factor can therefore supply a designated triangle. There are three
designated pairs per factor, and the nine pairs specify distinct faces.
\end{proof}

\section{Reproduction, provenance, and limits}
The complete exact certificate is
\path{results/fresh_compact29_three_triangle_star_constraint_20261018.json}.
Its producer, report, inherited data and full independent checker are
listed in \path{docs/compact29_triangle_constraints_map.md}. Reproduction
uses the producer's nonmutating \texttt{--verify} mode and the standalone
saved checker. The checker uses only the Python standard library and
imports no project module. It checks the finite certificate, not the
entire theory of hyperbolic polytopes from axioms.

The theorem is internally checked by exact computation and written
mathematical review. External refereeing, priority and full
proof-assistant verification are separate. No new compact dimension
bound follows at this stage: a quadrilateral can still consume the
current face-charge allowance. Improving the bound requires a further
charging or incidence theorem covering the remaining faces and all
relevant vertex types. No finite-volume construction is claimed.
\end{document}
